\section{Introducción}
\label{sec:introduction}

La crisis de deuda soberana en la Eurozona generó periodos de volatilidad extrema que desafiaron los modelos tradicionales de gestión de riesgos, convirtiendo a la prima de riesgo griega en un objeto de estudio paradigmático para comprender la inestabilidad financiera. Hasta ahora sabíamos que los mercados financieros en crisis no siguen un comportamiento normal, presentando "colas pesadas" y agrupamientos de volatilidad que sugieren que el riesgo no es constante en el tiempo. Por otro lado, la literatura ha establecido que la modelización de estos fenómenos puede abordarse desde la física estadística, asumiendo sistemas complejos con memoria, o desde la econometría financiera, asumiendo procesos estocásticos de varianza cambiante.

Sin embargo, existe una notable "laguna" científica en la intersección de estas disciplinas: la literatura tiende a tratar la Econofísica y la Econometría Financiera como compartimentos estancos, sin apenas estudios que comparen la eficacia explicativa y predictiva de ambos enfoques sobre una misma serie temporal de estrés extremo. No se ha publicado evidencia concluyente que enfrente directamente la capacidad de ajuste estructural de los modelos de difusión anómala contra la flexibilidad operativa de los modelos GARCH bajo un marco de inferencia Bayesiana. Esta investigación es relevante científicamente porque busca reconciliar ambas perspectivas, ofreciendo una visión dual de la dinámica del mercado.

El objetivo de este paper es cubrir dicho vacío comparando el desempeño del modelo de Difusión Anómala frente al modelo GARCH Bayesiano en la caracterización de la prima de riesgo griega. Partimos del supuesto teórico de que los mercados financieros no son completamente eficientes ni aleatorios (hipótesis de mercado eficiente), sino que poseen una estructura fractal con memoria a largo plazo que coexiste con ruido especulativo a corto plazo. Nuestro planteamiento teórico sintetizado sostiene que mientras el enfoque físico captura la naturaleza persistente ("tendencia") del riesgo sistémico, el enfoque econométrico es necesario para cuantificar la incertidumbre inmediata ("ruido") que enfrentan los agentes.